\documentclass[11pt]{amsart}

\usepackage[T1]{fontenc}
\usepackage{lmodern}
\usepackage{microtype}
\usepackage{mathtools,amssymb,amsthm}
\usepackage{array}
\usepackage[hidelinks]{hyperref}

\hypersetup{
  pdftitle={A Non-Standard (5,4) Broadcast of the Infinite Grid of Density 1/20}
}

\newtheorem{theorem}{Theorem}
\newtheorem{lemma}[theorem]{Lemma}
\newtheorem{proposition}[theorem]{Proposition}
\newtheorem{corollary}[theorem]{Corollary}
\theoremstyle{remark}
\newtheorem*{remark}{Remark}

\newcommand{\ZZ}{\mathbb{Z}}
\newcommand{\dens}{\delta}
\newcommand{\rec}{\operatorname{rec}}
\newcommand{\Tstd}{T}

\title[A non-standard $(5,4)$ broadcast of density $1/20$]
{A Non-Standard $(5,4)$ Broadcast of the Infinite Grid of Density $1/20$}

\date{}

\subjclass[2020]{05C69, 05B40, 94B25}
\keywords{broadcast domination, $(t,r)$ broadcast domination, infinite grid,
density, counterexample, perfect Lee code}

\begin{document}

\begin{abstract}
Blessing, Insko, Johnson, and Mauretour conjectured that the optimal $(t,r)$ and
$(t-1,r-2)$ broadcast dominating sets of $\ZZ\times\ZZ$ have equal density,
equivalently that $\dens_{t,r}=\dens_{t+1,r+2}$ for all $r\le t$. That
conjecture was already known to fail at other parameters; we add the instance
$(t,r)=(4,2)$, that is the pair of cells $(4,2)$ and $(5,4)$. The sublattice
$S=\langle(10,0),(4,2)\rangle$ of index $20$, obtained by doubling the perfect
Lee code of radius one, is a $(5,4)$ broadcast dominating set of $\ZZ^2$, so
\[
 \dens_{5,4}\le\tfrac1{20}<\tfrac1{18}=\dens_{4,2}.
\]
The certificate is the table of twenty coset receptions. $S$ is not a standard
broadcast: it has no vertex on any odd horizontal line, and
$\ZZ^2/S\cong\ZZ/2\oplus\ZZ/10$ is not cyclic. We do not determine
$\dens_{5,4}$; with a quoted lower bound of Shlomi it lies in
$\tfrac1{21}\le\dens_{5,4}\le\tfrac1{20}$.
\end{abstract}

\maketitle

\section{The conjecture}

Throughout, $\ZZ^2=\ZZ\times\ZZ$ is the infinite grid graph, whose graph
distance $d$ is the $\ell_1$ distance. Fix integers $1\le r\le t$. A vertex
$v$ of a set $S\subseteq\ZZ^2$ --- a \emph{tower} --- provides a vertex $u$ with
the signal $\max(0,t-d(u,v))$, so that a tower provides \emph{itself} with $t$.
The total signal received at $u$ is
\[
 \rec_S(u)\;=\;\sum_{v\in S}\max\bigl(0,\;t-d(u,v)\bigr),
\]
the sum including the term $v=u$ when $u\in S$. The set $S$ is a \emph{$(t,r)$
broadcast dominating set} if $\rec_S(u)\ge r$ for \emph{every} $u\in\ZZ^2$,
towers included. Its density is
$\dens(S)=\lim_{n\to\infty}\lvert S\cap[-n,n]^2\rvert/(2n+1)^2$, and
$\dens_{t,r}=\dens_{t,r}(\ZZ^2)$ denotes the least density of a $(t,r)$
broadcast dominating set of $\ZZ^2$. A \emph{standard} broadcast is a lattice
\[
 \Tstd(d,e)\;=\;\langle(d,0),(e,1)\rangle\;=\;\{(x,y)\in\ZZ^2:\ x\equiv ey\ (\mathrm{mod}\ d)\},
\]
one tower out of every $d$ on \emph{every} horizontal line; its density is $1/d$.

Blessing, Insko, Johnson, and Mauretour \cite{BIJM} close their paper with a
list of conjectures and open problems, whose fourth item reads, verbatim:

\begin{quote}
``We have seen that $(3,3)$ and $(2,1)$ broadcast dominating sets are equal for
large grids. We conjecture that the optimal $(t,r)$ and $(t-1,r-2)$ broadcast
dominating sets of $\ZZ\times\ZZ$ are equal. If this is true, then all $(t,r)$
broadcast domination problems for large grids can be reduced to $(t,2)$ and
$(t,1)$ broadcast domination problems.''
\end{quote}

\noindent
The parameter range $1\le r\le t$ is fixed earlier in that paper, and the
statement is made for the infinite grid, not for finite grids. Drews, Harris,
and Randolph restate it as a conjecture \cite{DHR}:
\emph{``For all $t,r\in\ZZ^+$, the optimal $(t,r)$ and $(t+1,r+2)$ broadcast
densities are equal.''} In the notation above the assertion is
\begin{equation}\label{eq:conj}
 \dens_{t,r}\;=\;\dens_{t+1,r+2}
 \qquad\text{for all } 1\le r\le t .
\end{equation}
We refute \eqref{eq:conj} at $(t,r)=(4,2)$.

\section{The witness}

\begin{theorem}\label{thm:main}
Let
\[
 S\;=\;\bigl\langle(10,0),\,(4,2)\bigr\rangle
   \;=\;\bigl\{(x,y)\in\ZZ^2:\ y\equiv0\ (\mathrm{mod}\ 2),\
        \ x\equiv2y\ (\mathrm{mod}\ 10)\bigr\}.
\]
Then $S$ is a $(5,4)$ broadcast dominating set of $\ZZ^2$ of density exactly
$1/20$. Consequently
\[
 \dens_{5,4}\;\le\;\tfrac1{20}\;<\;\tfrac1{18}\;=\;\dens_{4,2},
\]
so \eqref{eq:conj} is false at $(t,r)=(4,2)$.
\end{theorem}

Towers sit on even horizontal lines only; within an occupied line they are $10$
apart; each occupied line is shifted $4$ to the right of the occupied line two
below it. The two descriptions agree: $(x,y)=m(10,0)+n(4,2)$ has $y=2n$ even and
$x=10m+4n\equiv 2y\pmod{10}$, and conversely $y=2n$ together with
$x\equiv2y\pmod{10}$ gives $x-4n\equiv0\pmod{10}$. The basis matrix
$\begin{psmallmatrix}10&0\\4&2\end{psmallmatrix}$ has determinant $20$, so $S$
has index $20$ in $\ZZ^2$ and the twenty classes of $\ZZ^2/S$ are represented by
the pairs $(x,y)$ with $0\le x\le9$ and $y\in\{0,1\}$.

The certificate is the list of receptions at those twenty
representatives:
\begin{center}
\begin{tabular}{c|cccccccccc|c}
 & \multicolumn{10}{c|}{$x=0,1,\ldots,9$} & sum\\\hline
 $y=0$ & 5&4&5&4&5&4&5&4&5&4 & 45\\
 $y=1$ & 4&4&4&4&4&4&4&4&4&4 & 40\\\hline
 \multicolumn{11}{c}{} & 85
\end{tabular}
\end{center}
The minimum is $4=r$, attained at fifteen of the twenty classes. Two entries
worked out by hand: the shortest nonzero vector of $S$ is $(4,2)$, of
$\ell_1$ norm $6>4$, so a tower receives nothing but its own $5$, giving the
entry at $(0,0)$; and $(1,0)$ is at distance $1$ from the tower $(0,0)$ and at
distance $\ge5$ from every other tower --- the next nearest are $(4,2)$ and
$(2,-4)$, both at distance $5$ --- so it receives exactly $4$, which is the
tight case.

The sum of the twenty numbers is forced, which checks the table as a whole.

\begin{lemma}\label{lem:checksum}
Let $\Lambda\subseteq\ZZ^2$ be a sublattice of index $D$, let $t\ge1$, and let
$R$ be any set of coset representatives of $\ZZ^2/\Lambda$. Then
\[
 \sum_{u\in R}\rec_\Lambda(u)\;=\;W_t
 \;:=\;\sum_{z\in\ZZ^2}\max\bigl(0,t-\lVert z\rVert_1\bigr)
 \;=\;t+\sum_{m=1}^{t-1}4m(t-m),
\]
independently of $\Lambda$ and of $D$. In particular $W_5=5+16+24+24+16=85$.
\end{lemma}

\begin{proof}
Write $f(z)=\max(0,t-\lVert z\rVert_1)$, a nonnegative function of finite
support, so that $\rec_\Lambda(u)=\sum_{\lambda\in\Lambda}f(u-\lambda)$. As $u$
runs over $R$ and $\lambda$ over $\Lambda$, the difference $u-\lambda$ takes
each value of $\ZZ^2$ exactly once, because $R$ meets every coset of $\Lambda$
exactly once. Hence $\sum_{u\in R}\rec_\Lambda(u)=\sum_{z\in\ZZ^2}f(z)=W_t$.
Counting cells by $\ell_1$ norm --- one of norm $0$ and $4m$ of norm $m\ge1$ ---
gives the displayed formula.
\end{proof}

\begin{proof}[Proof of Theorem~\ref{thm:main}]
The twenty entries of the table sum to $45+40=85=W_5$, in agreement with
Lemma~\ref{lem:checksum}. That agreement checks the table jointly and does not
by itself bound the individual entries; read one by one, their minimum is $4$.
Since $\rec_S$ is invariant under
translation by $S$, and the twenty representatives meet every class of
$\ZZ^2/S$, we get $\rec_S(u)\ge4$ for every $u\in\ZZ^2$: the set $S$ is a
$(5,4)$ broadcast dominating set. It contains exactly one point per fundamental
domain, of area $20$, so $\dens(S)=1/20$.

For the consequence, $\dens_{4,2}=1/18$: Drews, Harris, and Randolph prove
$\dens_{t,2}=1/\bigl(2(t-1)^2\bigr)$ for all $t\ge2$ \cite{DHR}, which at $t=4$
gives $1/(2\cdot3^2)=1/18$. Since
$1/20<1/18$ and $\dens_{5,4}\le\dens(S)=1/20$, the two densities
$\dens_{4,2}$ and $\dens_{5,4}$ differ, contradicting~\eqref{eq:conj} at
$(t,r)=(4,2)$.
\end{proof}

\begin{remark}[$S$ as a doubled perfect Lee code]
Let $\varphi(X,Y)=(2X,2Y)$ and let
$L=\Tstd(5,2)=\langle(5,0),(2,1)\rangle=\{(X,Y):X\equiv2Y\ (\mathrm{mod}\ 5)\}$,
the perfect Lee code of radius $1$ \cite{GW} --- equivalently the optimal $(2,1)$
broadcast, of density $1/5$: every cell of $\ZZ^2$ is within distance $1$ of
exactly one codeword. Then $S=\varphi(L)$ and
$\dens(S)=\tfrac14\cdot\tfrac15=\tfrac1{20}$. Doubling doubles all $\ell_1$
distances, so at $t=5$ a tower $\varphi(v)$ reaches $\varphi(u)$ only when
$d(u,v)\le2$, with signal $5,3,1$ according as $d(u,v)=0,1,2$. On $2\ZZ^2$ this
settles the reception outright: a codeword receives its own $5$ and nothing else,
because the shortest nonzero vector of $L$ has norm $3$; a non-codeword receives
$3$ from the unique codeword at distance $1$ together with $1$ from each of the
two codewords at distance $2$, again $5$. The receptions off $2\ZZ^2$ are not
settled by this argument; the table above gives each of them as exactly~$4$.
\end{remark}

\begin{proposition}\label{prop:nonstandard}
$S$ is not a standard broadcast, nor a translate of one, nor the image of one
under any change of basis of $\ZZ^2$.
\end{proposition}

\begin{proof}
Every point of $S$ has even second coordinate, so $S$ misses the horizontal line
$y=1$ entirely, whereas every $\Tstd(d,e)$ contains a point on every horizontal
line; the same holds after translating. For the stronger statement, the basis
$\{(d,0),(e,1)\}$ of $\Tstd(d,e)$ has an entry equal to $1$, so the greatest
common divisor of the entries of its basis matrix is $1$, the invariant factors
are $(1,d)$, and $\ZZ^2/\Tstd(d,e)\cong\ZZ/d$ is cyclic. All four entries of the
basis matrix of $S$ are even, so its invariant factors are $(2,10)$ and
$\ZZ^2/S\cong\ZZ/2\oplus\ZZ/10$ is not cyclic. The isomorphism type of
$\ZZ^2/\Lambda$ is invariant under $GL_2(\ZZ)$, so no reflection, transposition,
or relabelling of a standard broadcast is $S$.
\end{proof}

\section{Scope}

What is refuted is the equality~\eqref{eq:conj} at the single instance
$(t,r)=(4,2)$. Each instance of \eqref{eq:conj} is a separate statement and needs
its own witness; we exhibit one only at $(4,2)$. The equality was already refuted
in print at other parameters: \cite{DHRcode} reports counterexamples, and
\cite{HvH} restates that Drews, Harris, and Randolph ``showed \dots\ that, in
fact, $\dens_{t+1,r+2}(\ZZ^2)<\dens_{t,r}(\ZZ^2)$ for several values of $t$ and
$r$''. The instance $(4,2)$ falls in the sub-case $r=2$ with $t$ even and $t>3$
singled out in \cite{DHRcode}.

\medskip\noindent
We do not determine $\dens_{5,4}$. Shlomi \cite{Shlomi} proves
$\dens_{t,r}\ge r/C_{t,r}$ with
$C_{t,r}=\tfrac23r^3-2r^2t+2rt^2+\tfrac13r$, which at $(t,r)=(5,4)$ gives
$C_{5,4}=84$ and hence
\[
 \tfrac1{21}\;=\;\tfrac4{84}\;\le\;\dens_{5,4}\;\le\;\tfrac1{20}.
\]
We made no search above index $22$, so nothing here excludes a $(5,4)$ broadcast
of density below $1/20$.

\section{Reproduction and verification}\label{sec:verif}

Everything above is determined by the two lines of data
\[
 \text{basis } \{(10,0),(4,2)\},
 \quad
 (x,y)\in S \iff y\equiv0\ (\mathrm{mod}\ 2),\ \ x\equiv2y\ (\mathrm{mod}\ 10),
\]
and a reader who wants only Theorem~\ref{thm:main} needs no computer: the twenty
receptions of the table are twenty sums of at most three terms each.

The program \texttt{verify.py} accompanying this note re-derives from those two
lines, in exact integer and rational arithmetic and using only the Python
standard library, the quantities printed above: the agreement of the predicate
with the span, the index $20$ and the transversality of the twenty
representatives, the $S$-periodicity of $\rec_S$, the two table rows entry by
entry with their sums and the value $W_5=85$, the minimum $4$, the density as an
exact rational, the shortest nonzero vector, the identity $S=2\,\Tstd(5,2)$, the
two arguments of Proposition~\ref{prop:nonstandard}, the value
$\dens_{4,2}=1/18$ from the quoted formula with the strict inequality
$1/20<1/18$, and the constant $C_{5,4}=84$ with $4/84=1/21<1/20$. Its $43$
checks all pass; the recorded run is \texttt{verify.output.txt}. That run also
contains checks of statements outside the scope of this note --- sublattice
censuses at indices $19$--$22$, standard broadcasts, further members of the
doubling family, and remarks on published tables and on a further question of
Drews, Harris, and Randolph. None of those is claimed here, and nothing above
depends on them.

The lower bounds of \cite{DHR} and \cite{Shlomi} are quoted, not reproved --- only
their arithmetic at the relevant parameters is checked, and the strict inequality
of Theorem~\ref{thm:main} rests on the theorem of \cite{DHR}. The program makes no
search above index $22$ and no literature search.

\begin{thebibliography}{9}

\bibitem{BIJM}
D.~Blessing, E.~Insko, K.~Johnson, and C.~Mauretour,
\emph{On $(t,r)$ broadcast domination numbers of grids},
Discrete Appl. Math. \textbf{187} (2015), 19--40.
\href{https://arxiv.org/abs/1401.2499}{arXiv:1401.2499}.
The conjecture refuted here is item~4 of the subsection ``Conjectures and Open
Problems''.

\bibitem{DHR}
B.~F. Drews, P.~E. Harris, and T.~W. Randolph,
\emph{Optimal $(t,r)$ broadcasts on the infinite grid},
Discrete Appl. Math. \textbf{255} (2019), 183--197.
\href{https://arxiv.org/abs/1711.11116}{arXiv:1711.11116}.
Restates the conjecture and proves $\dens_{t,2}=1/(2(t-1)^2)$.

\bibitem{DHRcode}
B.~F. Drews, P.~E. Harris, and T.~W. Randolph,
\emph{Computing upper bounds for optimal density of $(t,r)$ broadcasts on the
infinite grid},
\href{https://arxiv.org/abs/1712.00150}{arXiv:1712.00150}, 2017.
Carries the counterexamples at other parameters and the remark on $r=2$ with $t$
even and greater than $3$.

\bibitem{HvH}
R.~Herrman and P.~van Hintum,
\emph{$(t,r)$ broadcast domination in the infinite grid},
Discrete Appl. Math. \textbf{289} (2021), 270--280.
\href{https://arxiv.org/abs/1912.11560}{arXiv:1912.11560}.

\bibitem{Shlomi}
T.~Shlomi,
\emph{Bounds on $(t,r)$ broadcast domination of $n$-dimensional grids},
Discrete Math. Theor. Comput. Sci. \textbf{24} (2023), no.~1, paper
\texttt{dmtcs:5732}.
\href{https://doi.org/10.46298/dmtcs.5732}{doi:10.46298/dmtcs.5732};
\href{https://arxiv.org/abs/1908.07586}{arXiv:1908.07586}.
Lower bound $\dens_{t,r}\ge r/C_{t,r}$ with
$C_{t,r}=\tfrac23r^3-2r^2t+2rt^2+\tfrac13r$.

\bibitem{GW}
S.~W. Golomb and L.~R. Welch,
\emph{Perfect codes in the Hamming metric and the Lee metric},
SIAM J. Appl. Math. \textbf{18} (1970), 302--317.
The perfect Lee codes in $\ZZ^2$, of which $\Tstd(5,2)$ is the one of radius $1$.

\end{thebibliography}

\end{document}
